% Notas de aula — Capítulo 10: Forças Atmosféricas e Ventos
% Curso: Meteorologia Prática (baseado em R. Stull, Practical Meteorology, CC BY-NC-SA 4.0)
% Caixas verdes = conteúdo literal do quadro; linhas verdes = encenação.
\documentclass[11pt]{article}
\usepackage{../comum/preambulo}
\graphicspath{{../figuras/cap10/}}

\title{Capítulo 10 --- Forças Atmosféricas e Ventos\\
{\large Notas de aula (roteiro de quadro-negro)}}
\author{Curso de Meteorologia Prática}
\date{}

\begin{document}
\maketitle
\thispagestyle{fancy}

\noindent\textbf{Plano sugerido:} 2 aulas de 2\,h.
\emph{Aula 1:} \S\ref{sec:forcas}--\S\ref{sec:geo} (as forças, Coriolis,
vento geostrófico).
\emph{Aula 2:} \S\ref{sec:gradiente}--\S\ref{sec:continuidade}
(vento gradiente, atrito e camada limite, espiral de Ekman,
continuidade).

\section{A 2ª lei de Newton num planeta que gira}\label{sec:forcas}

Chegou a dinâmica. Os Caps.~1--9 montaram os campos ($P$, $T$, umidade)
e as cartas que os mostram \javimos{9}{análise de cartas}; agora
perguntamos: \emph{por que o vento sopra como sopra?} A resposta é a 2ª
lei de Newton por unidade de massa, escrita num referencial que gira.

\begin{noquadro}[a equação da aula --- montar termo a termo]
\[ \frac{D\vec{U}_h}{Dt} =
 \underbrace{-\frac{1}{\rho}\nabla_h P}_{\text{gradiente de pressão}}
 \;\underbrace{-\;f\,\hat{k}\times\vec{U}_h}_{\text{Coriolis}}
 \;\underbrace{-\;\frac{\vec U_h}{\tau_{atrito}}}_{\text{atrito (CL)}}
\qquad f = 2\Omega\sin\phi \]
HS: $f<0$ $\Rightarrow$ Coriolis desvia \textbf{para a esquerda}.
\end{noquadro}

\quadro{Apresentar cada força com seu desenho: (1) PGF apontando da alta
para a baixa, perpendicular às isóbaras; (2) Coriolis perpendicular ao
vento; (3) atrito contra o vento. Insistir: Coriolis não faz trabalho —
só gira.}

A força de gradiente de pressão é a única que \emph{cria} vento — as
outras duas o desviam ou o freiam. A PGF por unidade de massa é
$|\nabla P|/\rho$: com as isóbaras da carta do Cap.~9 e a régua de
111\,km/grau \javimos{1}{referenciais terrestres}, ela se mede
diretamente na carta.

\begin{formalivro}[parâmetro de Coriolis]
\[ f = 2\Omega\sin\phi,\qquad
\Omega = 7{,}292\times10^{-5}\ \mathrm{s^{-1}};
\quad f(-30°) = -7{,}3\times10^{-5}\ \mathrm{s^{-1}}. \]
\end{formalivro}

\begin{formadif}[de onde vem o $2\Omega\sin\phi$]
No referencial girante aparecem as forças fictícias centrífuga (já
absorvida no $g$ efetivo) e de Coriolis, $-2\vec\Omega\times\vec U$.
Projetando $\vec\Omega$ na vertical local: componente útil
$\Omega\sin\phi$; o fator 2 vem da soma dos efeitos ``mudança de raio''
e ``conservação de momento angular''. No equador $f = 0$ — o vento não
sente rotação horizontal ali (T2).
\end{formadif}

\begin{noquadro}[quem sente Coriolis?]
escala de tempo $1/|f| \sim 3$--4\,h\\
dura horas/dias (sinótico): sente \quad dura segundos/minutos (pia,
redemoinho): NÃO sente\\[2pt]
sem PGF, o vento faz \textbf{círculos inerciais}: período
$2\pi/|f|$, raio $V/|f|$
\end{noquadro}

Círculos inerciais são a primeira simulação do notebook (Caso~1) — e
aparecem de verdade em boias oceânicas à deriva. A hierarquia de escalas
espaciais e temporais (Fig.~10.24 do livro) diz quem entra no equilíbrio
em cada fenômeno: tornado (ciclostrófico), ciclone (gradiente), jato
(geostrófico) \veremos{11}{escalas da circulação geral}.

\section{Vento geostrófico}\label{sec:geo}

\quadro{Desenhar isóbaras retas e CONSTRUIR o equilíbrio em três atos:
1) ar parado sente só PGF e acelera para a baixa; 2) ganhando
velocidade, Coriolis gira; 3) estado final: PGF $\perp$ isóbaras,
Coriolis oposta, vento paralelo às isóbaras. Este desenho é o mais
importante do capítulo.}

\begin{formadif}[o equilíbrio central da meteorologia dinâmica]
Escoamento reto, sem atrito, estacionário: PGF $+$ Coriolis $= 0$:
\[ f\,U_g = -\frac{1}{\rho}\frac{\partial P}{\partial y},\qquad
   f\,V_g = +\frac{1}{\rho}\frac{\partial P}{\partial x}
\quad\Longrightarrow\quad
\vec U_g = \frac{1}{\rho f}\,\hat k\times\nabla_h P. \]
\end{formadif}

\begin{noquadro}[Buys-Ballot]
vento paralelo às isóbaras;\\
de costas para o vento no HS: baixa pressão à \textbf{direita}\\
(no HN: à esquerda)\\[2pt]
nas cartas de altitude: mesmo equilíbrio com alturas geopotenciais
\javimos{9}{cartas de 50\,kPa}
\end{noquadro}

\begin{exemplo}[régua sobre a carta]
Isóbaras de 4\,hPa espaçadas 300\,km, $\phi = -30°$,
$\rho = 1{,}2$\,kg/m$^3$:
\[ U_g = \frac{1}{\rho|f|}\frac{\Delta P}{\Delta n}
 = \frac{400}{1{,}2\times7{,}3\times10^{-5}\times3\times10^{5}}
 \simeq 15\ \mathrm{m/s}. \]
\emph{Verificação: vento sinótico típico; apertando para 150\,km dobra —
o ``apertado $=$ ventoso'' do Cap.~9, agora explicado.}
\end{exemplo}

Por que a natureza escolhe esse equilíbrio? Porque qualquer
desequilíbrio gera acelerações que Coriolis converte em giro até
alinhar — o ajustamento geostrófico. A rotação da Terra é o motivo de os
ventos \emph{contornarem} os centros de pressão em vez de preenchê-los
em horas; sem ela, não haveria ciclones durando dias
\veremos{13}{ciclones extratropicais}.

\section{Vento gradiente: curvas mudam tudo}\label{sec:gradiente}

\quadro{Desenhar baixa e alta com os triângulos de força; mostrar que a
resultante centrípeta exige vento diferente do geostrófico.}

Em torno de centros, sobra força centrípeta $V^2/R$:

\begin{formalivro}[vento gradiente]
\[ \frac{V^2}{R} =
\underbrace{\frac{1}{\rho}\left|\frac{\partial P}{\partial n}\right|}_{PGF}
 \mp |f|\,V
\quad\Rightarrow\quad
V_{baixa} < U_g < V_{alta}. \]
\end{formalivro}

\begin{noquadro}[baixas apertadas, altas moles]
baixas: Coriolis ajuda a curvar $\Rightarrow$ vento
\emph{sub}geostrófico; PGF pode ser gigante (furacões!)\\
altas: Coriolis briga $\Rightarrow$ vento \emph{super}geostrófico; e há
teto: $|\partial P/\partial n| \le \rho f^2R/4$ (T3)\\[2pt]
$\Rightarrow$ por isso baixas têm centros afiados e ventania; altas,
centros chatos e calmaria
\end{noquadro}

Esse contraste é visível em qualquer perfil de pressão de carta
(Fig.~10.14 do livro) — e explica por que só as \emph{baixas} podem
virar tempestades intensas \veremos{16}{ciclones tropicais}.

Caso extremo: $R$ pequeno e $V$ grande $\Rightarrow$ Coriolis
desprezível $\Rightarrow$ \textbf{equilíbrio ciclostrófico} (tornados,
redemoinhos — podem girar para qualquer lado!
\veremos{15}{tornados}). No extremo oposto (escoamentos rasos e
lentos, brisas fracas), o balanço PGF--atrito sem Coriolis é o vento
\textbf{antitríptico} — completa-se assim o cardápio do Stull:
geostrófico, gradiente, ciclostrófico, inercial e antitríptico.

\section{Atrito: a camada limite entra no jogo}\label{sec:atrito}

\quadro{Triângulo de forças com atrito: o vento cruza as isóbaras
$\sim$15--45° em direção à baixa. Desenhar a espiral convergindo para o
centro da baixa.}

Com atrito \veremos{18}{arrasto turbulento}, o equilíbrio de 3 forças
gira o vento em direção à baixa pressão.

\begin{noquadro}[a consequência sinótica mais importante]
atrito $\Rightarrow$ vento cruza isóbaras para a baixa\\
$\Rightarrow$ \textbf{convergência} em baixas: ar sobe $\to$ nuvem e
chuva\\
$\Rightarrow$ \textbf{divergência} em altas: ar desce $\to$ céu limpo\\[2pt]
a meteorologia inteira da carta de superfície está neste desenho
\end{noquadro}

É o elo que fecha com os Caps.~4--7: a subida forçada pela convergência
satura o ar \javimos{4}{LCL} e alimenta os nimbostratus
\javimos{6}{estratiformes}; e antecipa o ``bombeamento de Ekman'' que
acopla a camada limite aos ciclones \veremos{13}{ciclogênese}.

\subsection{Espiral de Ekman}

\begin{formadif}[a solução clássica]
Modelando o atrito como difusão turbulenta
$K\,\partial^2\vec U/\partial z^2$ e resolvendo o equilíbrio em função
de $z$:
\[ U(z) = U_g\left[1 - e^{-z/\delta}\cos(z/\delta)\right],\quad
   V(z) = U_g\,e^{-z/\delta}\sin(z/\delta),\quad
   \delta = \sqrt{2K/|f|}. \]
O vetor vento \emph{espirala} com a altura até o geostrófico em
$z \sim \pi\delta \sim 1$\,km. Verificação numérica no Caso~3; o T5
mostra que o transporte integrado é perpendicular a $\vec U_g$ — o
transporte de Ekman que também governa o oceano.
\end{formadif}

\section{Continuidade: o vento vertical nasce da convergência}
\label{sec:continuidade}

\quadro{Desenhar o cubo com fluxos entrando e saindo
\javimos{3}{o cubo euleriano}; se entra mais do que sai na horizontal,
sobe pelo topo.}

\begin{formalivro}[conservação de massa, forma incompressível]
\[ \frac{\partial u}{\partial x} + \frac{\partial v}{\partial y}
 + \frac{\partial w}{\partial z} = 0
\quad\Longrightarrow\quad
w(z_{topo}) = -\,z_{topo}\,
\overline{\left(\nabla_h\cdot\vec U_h\right)}. \]
\end{formalivro}

\begin{noquadro}[números do levantamento sinótico]
convergência de $10^{-5}$\,s$^{-1}$ em 1\,km $\Rightarrow$
$w \simeq 1$\,cm/s\\[2pt]
invisível — mas em poucas horas satura o ar (Cap.~4) e sustenta o
nimbostratus de uma frente
\end{noquadro}

Este é o quarto ingrediente da lista do Cap.~1 (continuidade
\javimos{1}{as quatro leis}) entrando em cena — e o levantamento que os
Caps.~12--13 usarão o tempo todo \veremos{12}{levantamento frontal}.
A \emph{cinemática} completa do campo de vento (divergência, que
acabamos de usar; vorticidade; deformação) será formalizada onde cada
peça trabalha: deformação na frontogênese \veremos{12}{deformação} e
vorticidade nos ciclones \veremos{13}{vorticidade}.

\textbf{Medindo ventos} (Stull \S10.10): anemômetro de conchas
(padrão das estações), de hélice, e o \textbf{sônico} — tempo de voo
do som em três eixos, rápido o bastante para medir turbulência
\veremos{18}{fluxos turbulentos}; em altitude, balão-piloto,
radiossonda GPS e os perfiladores \javimos{8}{sensoriamento}.

\section{Síntese da aula}

\begin{noquadro}[mapa final --- canto do quadro]
\[ f = 2\Omega\sin\phi \;\to\;
   \vec U_g = \frac{1}{\rho f}\hat k\times\nabla P \;\to\;
   \frac{V^2}{R}\ \text{(gradiente)} \;\to\;
   \text{atrito}\Rightarrow\text{convergência} \;\to\; w \]
do equilíbrio geostrófico à chuva da baixa pressão: quatro passos
\end{noquadro}

\section{Exercícios complementares}\label{sec:exercicios}

Soluções (professor) em \texttt{cap10\_solucoes.ipynb}.

\subsection*{Teóricos}

\begin{enumerate}[label=\textbf{T\arabic*.},itemsep=4pt]
  \item Mostre que, sem PGF e atrito, as equações do movimento dão
        círculos inerciais de período $2\pi/|f|$ e raio $V/|f|$; calcule
        ambos para $V = 10$\,m/s em $\phi = -23{,}5°$ e $-60°$.
  \item Derive o vento geostrófico e mostre que ele \emph{não pode}
        existir no equador; o que equilibra a PGF lá?
  \item Resolva a quadrática do vento gradiente para baixas e altas;
        mostre que em altas só há solução real se
        $|\partial P/\partial n| \le \rho f^2 R/4$ — o ``teto'' de PGF
        das altas.
  \item A partir do equilíbrio de 3 forças com atrito linear
        ($-\vec U/\tau$), deduza o ângulo de cruzamento das isóbaras
        $\tan\alpha = 1/(f\tau)$ e avalie para $\tau = 0{,}5$ e 3\,h.
  \item Verifique por substituição que a espiral de Ekman satisfaz o
        sistema $-fV = K U''$, $f(U - U_g) = K V''$, e mostre que o
        transporte integrado na camada (transporte de Ekman) é
        perpendicular ao vento geostrófico.
\end{enumerate}

\subsection*{Numéricos (no notebook)}

\begin{enumerate}[label=\textbf{N\arabic*.},itemsep=4pt]
  \item Integre trajetórias inerciais em $\phi = -10°, -30°, -60°$ e
        confirme período e raio do T1; adicione um vento médio e observe
        as cicloides.
  \item Para o campo de pressão sintético do Caso~2, compare o vento
        geostrófico com o vento gradiente ao redor da baixa e
        quantifique o quanto o gradiente é subgeostrófico no raio de
        máximo aperto.
  \item Resolva a camada de Ekman numericamente (diferenças finitas)
        para $K$ variável com a altura e compare com a espiral analítica
        de $K$ constante.
  \item Com o vento do Caso~2 mais um atrito de 20° de cruzamento,
        calcule a divergência horizontal e o $w$ no topo da camada
        limite; mapeie onde o ar sobe.
\end{enumerate}

\vspace{1em}
\noindent\rule{\linewidth}{0.4pt}\\
{\scriptsize Material derivado de R.~Stull, \emph{Practical Meteorology}
(CC BY-NC-SA 4.0); este documento herda a mesma licença.}

\end{document}
